\section{Conclusion}

We propose the sequential meta-analysis research trace (SMART). We take a sequence of studies, each one targeting the same parameter, and containing an estimate and a standard error. Applying a Bayesian approach, we plot sequentially updated 
posterior estimates of the target parameter and calculate each study's research contribution,
defined as the change in
its posterior relative to the prior literature. 
Our approach complements classical methods for meta-analysis.
Classical meta-analysis analyzes all the studies simultaneously to
produce a blended estimate of the target parameter, as well as weights that describe each study's
retrospective importance. Here we analyze the studies sequentially,
which leads to the same blended final estimate but assesses the importance
of each study at the time it appeared.


We are by no means the first to consider or conduct sequential meta-analysis; other work also updates a Bayesian meta-analysis model one study at a time \citep{Higgins_Whitehead_Simmonds_2011,Van_Der_Tweel_Bollen_2010,Bollen_Uiterwaal_Van_Vught_Van_Der_Tweel_2006}. But we additionally quantify the change in posterior, proposing that meta-analysts use the Wasserstein distance for gauging the importance of a study at the time it appeared.  This yardstick for quantifying research contributions can recognize the substantial impact of small studies that move a literature in a distinctive direction --- studies that would ordinarily be accorded little weight by a random effects meta-analysis, whether sequential or retrospective.  Although our approach to meta-analysis does not presuppose a specific meta-analytic model, we believe that models that take note of methodological ``labels’’ may offer insights that might otherwise be missed when summarizing literatures with active methodological debates.  


Our simulations and empirical example show that methodological innovations are characterized very differently by models that explicitly acknowledge the possibility that new insights may impeach the evidence that came before.
The profound influence of the Card and Krueger minimum wage study, for example, reflects its methodological insights rather than the precision of its empirical estimate.  In the social sciences, methodological upheaval is not simply an abstract possibility; one cannot appreciate the development of influential research literatures without taking notice of the sharp break that occurred due to the advent of natural experiments \citep{Ehrenberg_92}, regression discontinuity designs \citep{Thistlethwaite_60,Imbens_08}, or randomized control trials \citep{Angrist_10,Banerjee_16}.

We emphasize that SMART is a way of looking at a research literature.
Like all model-based methods, it comes with assumptions and
simplifications. We assume the parameter of interest is fixed over
time and that the sequence of studies we encounter constitute independent random draws from a sampling distribution.  We make the same parametric assumptions about study
likelihoods and priors that are posited in the classical
random-effects meta-analysis model, but when applying a ``labeled’’ meta-analysis model, we further assume that we can quantify the relative uncertainty over a set of 
methodological biases. 

That said, the model that underlies SMART provides potential insights
into the trajectory of a scientific literature. It can detect when a
literature loses momentum as it converges on a consensus answer with
increasing precision. Or it can show that a literature has been
revitalized by methodological innovations that expose the deficiencies
of prior approaches. With an extension to a dynamic model, SMART could further detect a changing scientific landscape--where the underlying estimand changes, be it due to temporal effects or a change in the scientific question itself.
